\documentclass[9pt,twocolumn,letterpaper]{article}
\usepackage{visstyle}
\usepackage{pgfplots}
\pgfplotsset{compat=1.18}

% Generated by scripts/make_figures.py -- do not edit.
\newcommand{\snapshot}{8 August 2026}
\newcommand{\vlhuman}{8{,}020}
\newcommand{\vlauthors}{236}
\newcommand{\vlpeak}{1{,}855}
\newcommand{\vlpeakyear}{2015}
\newcommand{\vlrecent}{55}
\newcommand{\vldecline}{97}
\newcommand{\vlbotshare}{58}
\newcommand{\vlbusfactor}{2}
\newcommand{\ecorepos}{14}
\newcommand{\ecocommits}{26{,}902}
\newcommand{\ecohumans}{666}
\newcommand{\ecobusfifty}{3}
\newcommand{\ecobuseighty}{8}
\newcommand{\ecotopthree}{66}
\newcommand{\ecorecent}{483}
\newcommand{\ecorecenthumans}{62}
\newcommand{\dthreecommits}{7{,}681}
\newcommand{\dthreehumans}{252}
\newcommand{\dthreetop}{77}
\newcommand{\dthreerecent}{2}
\newcommand{\schemareleases}{118}
\newcommand{\schemafirstdefs}{170}
\newcommand{\schemafirsttag}{v2.0.0}
\newcommand{\schemafirstyear}{2017}
\newcommand{\schemapeakdefs}{458}
\newcommand{\schemafirstkb}{307}
\newcommand{\schemalastmb}{1.79}
\newcommand{\schemafrozen}{6}
\newcommand{\schemafrozenfrom}{v6.2.0}
\newcommand{\schemafrozenfromdate}{2025-06-27}
\newcommand{\schemafrozento}{v6.4.3}
\newcommand{\ntimeline}{10}
\newcommand{\ndormant}{5}
\newcommand{\falconlast}{2025-07-30}
\newcommand{\voyagerlast}{2021-10-14}
\newcommand{\dracolast}{2023-08-07}
\newcommand{\lyralast}{2021-04-02}
\newcommand{\protovislast}{2013-11-05}
\newcommand{\mosaicrecent}{180}
\newcommand{\mplrecent}{1{,}194}
\newcommand{\mplbus}{7}


\vistitle{Who Maintains the Grammar?\\Bots, Bus Factors, and the Afterlife of Visualization Research Software}
\visauthor{Anonymous Submission}
\visaffil{}

\begin{document}
\maketitle

\begin{abstract}
The visualization community increasingly ships software with its papers, and
some of that software has become infrastructure that the rest of the field
builds on. We know a good deal about whether research code is \emph{released};
we know very little about whether it is \emph{maintained}. We mine the full
commit histories of \ecorepos{} repositories in the Protovis\,$\rightarrow$\,D3\,%
$\rightarrow$\,Vega\,$\rightarrow$\,Vega-Lite\,$\rightarrow$\,Mosaic lineage,
together with ten comparison toolkits, and report four measurements. Effort is
extraordinarily concentrated: \ecobusfifty{} people account for half of
\ecocommits{} human commits by \ecohumans{} people. Human effort collapses
rather than plateaus: Vega-Lite fell from \vlpeak{} human commits in
\vlpeakyear{} to \vlrecent{} in the last twelve months, a \vldecline\% decline.
Automation fills the vacuum: \vlbotshare\% of Vega-Lite's recent commits are
written by bots, and in three sibling repositories bots write more than four
commits in five. And award-winning artifacts go dormant early --- Voyager's
last human commit predates its Test-of-Time award by four years. We argue that
release-time reproducibility checks measure the wrong moment, and sketch a
lightweight maintenance disclosure that would make the situation legible.
\end{abstract}

\begin{indexterms}
Visualization software, research software engineering, repository mining,
reproducibility, sustainability, grammars of graphics.
\end{indexterms}

\section{Introduction}

A visualization paper that ships code is now the norm rather than the
exception, and the community has built real machinery to encourage it: the
Graphics Replicability Stamp Initiative, IEEE VIS Open Practices, and a
growing literature on whether artifacts are available and whether they
run~\cite{isenberg2024stamps}. That machinery evaluates a single instant --- the
moment of publication. It asks whether the code was released and whether a
reviewer could execute it.

For most papers that is the right question, because most research prototypes
are arguments rather than products. But a handful of systems escape that
category. Vega-Lite~\cite{satyanarayan2017vegalite} is not merely a paper with
an artifact attached; it is the substrate that a substantial part of the
field's subsequent empirical work runs on, from natural-language charting
interfaces to notebook toolchains. When a research system becomes
infrastructure, availability at publication time stops being the interesting
property. The interesting property is whether anyone is still there.

Nobody has measured that. This paper does, for the best-documented lineage the
visualization community has: five generations of published systems from one
research group, each with its own paper at a visualization venue, spanning
seventeen years of continuous version control. We mine the commit histories of
\ecorepos{} repositories in that lineage plus ten widely-used comparison
toolkits, and ask who is actually writing the commits.

The answer is uncomfortable, and it is not simply a story of software aging.
Our contributions are: (1) the first longitudinal measurement of maintenance
effort across a visualization research lineage; (2) four principal findings --- extreme
effort concentration, collapse rather than plateau of human contribution, bot
substitution, and early dormancy of award-winning artifacts --- each quantified
against comparison toolkits that show the pattern is contingent, not
inevitable; and (3) a concrete, cheap proposal for making maintenance status
visible at the moment readers actually need it.

\section{Related Work}

\subsubsection{Artifact availability}
Isenberg's study of reproducibility stamps~\cite{isenberg2024stamps}
characterizes how visualization papers release code and how that practice has
spread. It is the closest neighbour to this work and, we think, exactly half of
the picture: it measures release, we measure survival. The community has
precedent for turning its own output into
data~\cite{isenberg2017vispubdata,chen2021vis30k}, but it has aimed that
instrument at papers and figures rather than at code.

\subsubsection{Mining visualization artifacts}
Prior work has mined visualization \emph{artifacts and usage}: corpora of D3
visualizations crawled from the web~\cite{hoque2020searching}, Stack Overflow
questions about D3, and studies of how Observable users decompose code. These
treat the toolkit as a fixed backdrop and study what people make with it. We
invert this and study the toolkit's own repository.

\subsubsection{Sustainability in software engineering}
Bus factor and contributor concentration are standard instruments in software
engineering, as are studies of why open-source projects fail~%
\cite{coelho2017fail,avelino2016truck,hassan2008road}. Concentration around one
or two maintainers is unremarkable for open source in general. Our claim is not
that the Vega ecosystem is unusual \emph{as open source}; it is that a field
which has made this software load-bearing has never checked, and that the
numbers are worse than the field's reliance on them would justify.

\subsubsection{The research--practice gap}
The VisGap workshop has convened at EuroVis annually since 2022 to discuss the
distance between visualization research and visualization software. That
discussion has been conducted almost entirely through position papers and
experience reports. We offer it a quantitative baseline.

\section{Corpus and Method}

We study \ecorepos{} repositories from a single research lineage ---
Protovis~\cite{bostock2009protovis}, D3~\cite{bostock2011d3},
Vega~\cite{satyanarayan2016reactivevega},
Vega-Lite~\cite{satyanarayan2017vegalite} and Mosaic~\cite{heer2024mosaic},
plus the prototypes and support packages published alongside them
(Lyra, Voyager~\cite{wongsuphasawat2016voyager}, Draco~\cite{moritz2019draco},
Falcon, Altair~\cite{vanderplas2018altair}, and the editor, embed, tooltip and
datasets packages) --- and ten comparison toolkits
(Matplotlib~\cite{hunter2007matplotlib}, Bokeh, plotly.py, Chart.js,
ECharts~\cite{li2018echarts}, deck.gl, ggplot2~\cite{wickham2010layered},
Observable Plot, Gephi, and the ten \repo{d3-*} core modules analysed as one
family). The lineage was chosen because it is the only case we know of in which
every generation has a paper at a visualization venue, which makes research
provenance a controlled property rather than a judgement call. It is also a
lineage with an unusually explicit intellectual debt: each generation refines a
grammar of graphics~\cite{wilkinson2005grammar} into executable form, which is
part of why its artifacts became infrastructure rather than demonstrations.

Each repository is cloned blobless and its full history walked with
\repo{git log --no-merges}. Merge commits are excluded as bookkeeping. Author
names are lower-cased and passed through a small, explicit alias table that
merges handle/real-name pairs for the same person; this is the only manual
judgement in the pipeline, and it is conservative in the direction that
\emph{reduces} apparent concentration. Bot commits are identified by author
name and reported separately rather than discarded, because their share is
itself a result. All windows are measured against a pinned snapshot date of
\snapshot{}, so re-running the analysis reproduces the numbers below rather
than drifting.

Commit counts are a crude proxy for effort, and we lean on them only for
order-of-magnitude claims and for changes over time within a single repository,
where the proxy's biases are at least stable. Section~\ref{sec:threats}
revisits this.

\begin{figure}[t]
\centering
% Generated by scripts/make_figures.py -- do not edit.
\begin{tikzpicture}
\begin{axis}[
  width=\columnwidth, height=0.60\columnwidth,
  ybar stacked, bar width=6pt,
  xmin=2013, xmax=2027,
  xtick={2014,2016,2018,2020,2022,2024,2026},
  x tick label style={font=\scriptsize, rotate=45, anchor=east,
                       /pgf/number format/1000 sep={}},
  ylabel={Commits}, ymin=0,
  label style={font=\scriptsize}, tick label style={font=\scriptsize},
  legend style={font=\scriptsize, at={(0.98,0.98)}, anchor=north east,
                 draw=none, fill=white, fill opacity=0.85, text opacity=1},
  grid=major, grid style={gray!18},
  axis lines*=left,
]
\addplot[fill=vislink, draw=none] coordinates {(2014,202) (2015,1855) (2016,1387) (2017,1739) (2018,898) (2019,701) (2020,648) (2021,192) (2022,85) (2023,132) (2024,72) (2025,71) (2026,38)};
\addlegendentry{human}
\addplot[fill=black!22, draw=none] coordinates {(2014,0) (2015,0) (2016,0) (2017,0) (2018,0) (2019,116) (2020,267) (2021,328) (2022,400) (2023,231) (2024,63) (2025,83) (2026,45)};
\addlegendentry{bot}
\end{axis}
\end{tikzpicture}

\caption{Commits to \repo{vega/vega-lite} per year, separated into human and
automated authorship. Human contribution peaks in \vlpeakyear{} and declines by
\vldecline\% to \vlrecent{} commits in the last twelve months; bots, absent
before 2019, now write \vlbotshare\% of the repository's commits. The final bar
covers a partial year, to the \snapshot{} snapshot.}
\label{fig:bots}
\end{figure}

\section{Findings}

\subsection{Half the ecosystem is three people}

Across the \ecorepos{} lineage repositories there are \ecocommits{} human
commits from \ecohumans{} distinct people. Half of that work was done by
\ecobusfifty{} of them and eighty percent by \ecobuseighty{}; the top three
contributors alone write \ecotopthree\% of all commits, and one person has
committed to all \ecorepos{} repositories. In the last twelve months the whole
ecosystem received \ecorecent{} human commits from \ecorecenthumans{} people.

The D3 family is more concentrated still: of \dthreecommits{} human commits a
single author wrote \dthreetop\%, and the last twelve months brought
\dthreerecent{} human commits across all eleven repositories combined. Ten of
the \repo{d3-*} core modules share a last human commit in 2023 or 2024.

\subsection{Contribution collapses, it does not plateau}

Figure~\ref{fig:bots} shows the shape of the decline. If Vega-Lite were simply
a finished piece of software we would expect contribution to fall to a low,
stable floor. Instead it falls by \vldecline\% from its \vlpeakyear{} peak and
keeps falling, while the project's own backlog grows: at the snapshot date
Vega-Lite had 718 open and 2{,}872 closed issues, a 20\% open ratio, against
D3's 5 open issues.\footnote{Issue counts come from the GitHub API on the
snapshot date and, unlike everything else reported here, are not reproducible
from the git history; they are recorded in \repo{data/issues.json}.} D3 can
reasonably be described as finished. A project with a fifth of its issues
outstanding and two maintainers cannot.

\subsection{Bots have replaced people}

Before 2019 no commit in Vega-Lite was written by a bot. Bots reached 63\% of
commits in 2021, and over the last twelve months they account for
\vlbotshare\%. The pattern is sharper in the support packages that downstream
users actually depend on: bots write 95.9\% of recent commits to
\repo{vega-tooltip}, 87.8\% to \repo{vega-embed}, and 81.1\% to
\repo{vega-datasets}. In three repositories the all-time leading committer is
now Dependabot.

This is not nothing --- dependency bumps are real maintenance, and a project
whose bots keep running is in better shape than one whose bots have stopped.
But it is maintenance that cannot triage an issue, review a design proposal, or
decide what the grammar should mean.

\subsection{The grammar stopped growing}

Vega-Lite ships a machine-readable JSON Schema with every release, which makes
its expressiveness measurable rather than arguable --- a property no other
visualization grammar we know of provides. Across \schemareleases{} releases
the schema grew from \schemafirstdefs{} definitions and \schemafirstkb\,KB at
\schemafirsttag{} (\schemafirstyear{}, the release accompanying the paper) to
\schemapeakdefs{} definitions and \schemalastmb\,MB --- and then stopped. The
schema has been byte-identical across the last \schemafrozen{} releases, from
\schemafrozenfrom{} (\schemafrozenfromdate{}) to \schemafrozento{}. Whatever the
project is now doing, it is not extending the grammar.

\begin{figure}[t]
\centering
% Generated by scripts/make_figures.py -- do not edit.
\begin{tikzpicture}
\begin{axis}[
  width=\columnwidth, height=0.78\columnwidth,
  xmin=2008.5, xmax=2027.10,
  ymin=-0.8, ymax=9.8,
  ytick={0,1,2,3,4,5,6,7,8,9},
  yticklabels={{Protovis\,\tiny InfoVis'09},{D3\,\tiny InfoVis'11},{Lyra\,\tiny EuroVis'14},{Vega\,\tiny InfoVis'16},{Vega-Lite\,\tiny InfoVis'17},{Voyager\,\tiny InfoVis'15},{Altair\,\tiny JOSS'18},{Falcon\,\tiny CHI'19},{Draco\,\tiny InfoVis'18},{Mosaic\,\tiny VIS'24}},
  xtick={2010,2014,2018,2022,2026},
  label style={font=\scriptsize}, tick label style={font=\scriptsize},
  x tick label style={font=\scriptsize, /pgf/number format/1000 sep={}},
  y tick label style={font=\scriptsize},
  xlabel={},
  grid=none,
  axis lines*=left, y axis line style={draw=none}, ytick style={draw=none},
  clip=false,
  legend style={font=\scriptsize, at={(0.5,-0.14)}, anchor=north,
                 legend columns=2, draw=none, fill=none},
]
\draw[line width=3.2pt, black!20] (axis cs:2009.14,0) -- (axis cs:2013.84,0);
\draw[line width=3.2pt, black!20] (axis cs:2010.74,1) -- (axis cs:2026.40,1);
\draw[line width=3.2pt, black!20] (axis cs:2012.78,2) -- (axis cs:2021.25,2);
\draw[line width=3.2pt, black!20] (axis cs:2013.09,3) -- (axis cs:2026.57,3);
\draw[line width=3.2pt, black!20] (axis cs:2014.88,4) -- (axis cs:2026.58,4);
\draw[line width=3.2pt, black!20] (axis cs:2015.05,5) -- (axis cs:2021.78,5);
\draw[line width=3.2pt, black!20] (axis cs:2015.71,6) -- (axis cs:2026.57,6);
\draw[line width=3.2pt, black!20] (axis cs:2016.76,7) -- (axis cs:2025.57,7);
\draw[line width=3.2pt, black!20] (axis cs:2017.41,8) -- (axis cs:2023.60,8);
\draw[line width=3.2pt, black!20] (axis cs:2022.93,9) -- (axis cs:2026.60,9);
\addplot[only marks, mark=*, mark size=2.2pt, color=black!45]
  coordinates {(2013.84,0) (2021.25,2) (2021.78,5) (2025.57,7) (2023.60,8)};
\addlegendentry{no human commit in 12 months}
\addplot[only marks, mark=*, mark size=2.2pt, color=vislink]
  coordinates {(2026.40,1) (2026.57,3) (2026.58,4) (2026.57,6) (2026.60,9)};
\addlegendentry{still maintained}
\draw[dashed, black!35] (axis cs:2026.60,-0.8) --
                        (axis cs:2026.60,9.6);
\end{axis}
\end{tikzpicture}

\caption{Active lifetime of each published system in the lineage, from first
commit to last commit by a human. Systems with no human commit in the twelve
months before the snapshot are marked in grey. \ndormant{} of the
\ntimeline{} are dormant by that measure, among them a best-paper winner
(Draco) and a test-of-time winner (Voyager).}
\label{fig:timeline}
\end{figure}

\subsection{Awards arrive after the lights go out}

Figure~\ref{fig:timeline} places each published system on a timeline from its
first commit to its last human commit. Of the \ntimeline{}, \ndormant{} have
received no human commit in the last twelve months: Protovis since
\protovislast{}, Lyra since \lyralast{}, Draco --- a best-paper winner --- since
\dracolast{}, and Falcon since \falconlast{}.

The sharpest case is Voyager~\cite{wongsuphasawat2016voyager}, which received
the InfoVis Test-of-Time award in 2025. Its last human commit was
\voyagerlast{}: the artifact had been dormant for four years at the moment the
community honoured its durability. The award was for the ideas, and the ideas
have indeed lasted. But the gap between what we celebrate and what we sustain
is precisely the gap this paper is about.

\subsection{The pattern is contingent, not inevitable}

Table~\ref{tab:maintenance} makes clear that this is not simply what happens to
ten-year-old visualization software. Matplotlib, older still and far larger,
received \mplrecent{} human commits in the same twelve months with a bus factor
of \mplbus{}. ECharts and Bokeh, both foundation- or company-governed,
remain healthily active. What distinguishes them is not age or language but
governance: they have an institutional home that outlives any individual
contributor, and research-group software generally does not.

Effort in the lineage has not disappeared so much as migrated. Mosaic, the
newest generation, received \mosaicrecent{} human commits in the last twelve
months --- more than Vega-Lite and Vega combined. The lab is still producing
excellent research software. It is producing the \emph{next} one.

\begin{table}[t]
\centering
\caption{Maintenance profile at \snapshot{}. \emph{Bus\,50} is the number of
authors accounting for half of all human commits; \emph{Top} is the leading
author's share; \emph{12mo} counts human commits in the last twelve months and
\emph{Bot} the share of commits in that window written by bots.}
\label{tab:maintenance}
\small
\begin{tabular}{@{}lrrrrrr@{}}
\toprule
Repository & Commits & Authors & Bus\,50 & Top & 12mo & Bot \\
\midrule
% Generated by scripts/make_figures.py -- do not edit.
vega-lite & 8{,}020 & 236 & 2 & 48\% & 55 & 58\% \\
vega & 5{,}980 & 165 & 1 & 66\% & 80 & 49\% \\
d3 & 3{,}822 & 154 & 1 & 68\% & 2 & 0\% \\
altair & 3{,}082 & 204 & 3 & 42\% & 117 & 11\% \\
mosaic & 1{,}028 & 34 & 1 & 70\% & 180 & 25\% \\
matplotlib & 37{,}691 & 1{,}976 & 7 & 11\% & 1{,}194 & 3\% \\
bokeh & 17{,}788 & 734 & 3 & 30\% & 280 & 0\% \\
echarts & 8{,}355 & 351 & 4 & 17\% & 193 & 8\% \\
ggplot2 & 5{,}247 & 384 & 2 & 43\% & 92 & 0\% \\
plot & 2{,}125 & 28 & 1 & 80\% & 37 & 0\% \\
\bottomrule

\end{tabular}
\end{table}

\section{Discussion}

\subsubsection{We are measuring the wrong moment}
Reproducibility stamps and artifact-availability checks certify a system at
publication, when enthusiasm and funding are at their peak and the student who
wrote it has not yet graduated. Every system in our corpus would have passed
such a check. The information a reader needs in 2026 --- is this maintained, by
whom, and should I build on it --- is exactly what the badge does not carry.

\subsubsection{A maintenance disclosure}
We propose something deliberately small, because proposals that require
sustained volunteer labour are self-defeating given the finding above. Every
number in this paper was computed from public git history by a script that runs
in minutes. A venue could attach to each artifact a generated, dated line
reporting human commits in the last twelve months, distinct recent authors, bot
share, and date of last human commit --- and, critically, recompute it on the
digital library page rather than freezing it at publication. This is not a
quality judgement and should not be read as one. Dormant is not the same as
bad: Protovis is dormant because D3 replaced it, which is a success. The point
is to make the distinction between \emph{finished}, \emph{superseded} and
\emph{abandoned} visible instead of leaving each reader to guess.

\subsubsection{An honest tension}
Researchers are rewarded for the next paper, not for the previous system's
issue tracker, and we do not think individual authors are at fault here --- the
migration of effort to Mosaic is exactly what the incentive structure asks for.
If the community wants its infrastructure maintained, the realistic routes are
institutional homes for the few systems that become load-bearing, or explicit
acknowledgement that research software has a half-life and planning for
succession rather than pretending otherwise.

\section{Threats to Validity}
\label{sec:threats}

Commit counts measure activity, not value; a single well-judged commit can be
worth a hundred. We use them for large ratios and within-repository trends
rather than cross-project quality comparisons, but the proxy remains crude, and
work that never reaches a commit --- issue triage, design review, user support
--- is invisible to it. Our alias table merges identities we could confirm and
will have missed others, which biases concentration \emph{downward}, making our
central finding conservative. Bot classification by author name will
misclassify humans committing through automation. The twelve-month window is a
choice; a longer window softens the picture, a shorter one sharpens it.
Finally, one lineage is not the field --- it is the best-documented case, not a
random sample, and the comparison toolkits are a convenience sample of widely
used projects rather than a systematic one. We report a pattern within a
well-characterized corpus and do not claim it generalizes to visualization
research software as a whole.

\section{Conclusion}

We measured who maintains the visualization community's most-used research
software and found \ecobusfifty{} people behind half of it, human contribution
down \vldecline\% in its flagship project, bots writing the majority of recent
commits, a grammar that has stopped growing, and a test-of-time award given to
an artifact that had been dormant for four years. Comparison with
foundation-governed toolkits of the same age shows none of this is a law of
software aging.

Our field has taught itself to release code. The measurements here suggest the
harder and now more urgent question is what happens in the decade afterwards.

{\small
\bibliographystyle{abbrv}
\bibliography{refs}
}

\end{document}
